\pdfminorversion=7%
\documentclass[aspectratio=169,mathserif,notheorems]{beamer}%
%
\xdef\bookbaseDir{../../bookbase}%
\xdef\sharedDir{../../shared}%
\RequirePackage{\bookbaseDir/styles/slides}%
\RequirePackage{\sharedDir/styles/styles}%
\toggleToGerman%
%
\subtitle{8.~Fabrik-Datenbank:~Tabelle \emph{product}}%
%
\begin{document}%
%
\startPresentation%
%
\section{Einleitung}%
%
\begin{frame}%
\frametitle{Einleitung}%
\begin{itemize}%
\item Nun wollen wir die eigentliche Datenbank entwickeln.%
\item<2-> Normalerweise würden wir dabei einem Designprozess folgen, wir würden die Anforderungen aufschreiben und analysieren, ein konzeptuelles Modell entwickeln und daraus ein logisches Modell ableiten\dots%
\item<3-> {\dots}das machen wir jetzt nicht -- unser Ziel ist \emph{Learning by Doing.}%
\end{itemize}%
\end{frame}%
%
\begin{frame}%
\frametitle{Wir nutzen eine Relationale Datenbank}%
\begin{itemize}%
\item Wahrscheinlich kennen Sie bereits Tabellenkalkulationssoftware wie \microsoftExcel\cite{P2020MSS2ABG,A2024TSAFMSS2,W2018MSSDB} oder \libreofficeCalc\cite{DF2024LTDF,GL2012LTSOOSSCBAFACSOL,S2022L7PFEUU}.%
\item<2-> Auch dort sind Daten in Tabellen organisiert.%
\item<3-> In einer relationalen Datenbank sind die Spalten der Tabellen aber \emph{typisiert}\uncover<4->{: Sie können Text eingeben, der dem vorgegebenen Datentyp entspricht.}%
\item<5-> Und es kann mehrere verschiedene Tabellen geben.%
\item<6-> Und die Datensätze in verschiedenen Tabellen können miteinander fest verbunden sein.%
\end{itemize}%
\end{frame}%
%
\begin{frame}%
\frametitle{Struktur der Datenbank}%
\begin{itemize}%
\item Dieses Format erlaubt es uns, unsere Daten nach Kategorie einzuteilen.%
\item<2-> Wir werden eine Tabelle für die \emph{Producte} unseres Betriebs erstellen.%
\item<3-> Wir werden eine Tabelle für die \emph{Kunden} unseres Betriebs erstellen.%
\item<4-> Wir werden eine Tabelle für die \emph{Bestellungen} der Kunden unseres Betriebs erstellen.%
\item<5-> Fangen wir mit der Tabelle für die Produkte an.%
\end{itemize}%
\end{frame}%
%
\section{Erstellen der Tabelle}%
%
\begin{frame}%
\frametitle{Name der Tabelle}%
\begin{itemize}%
\item For many programming languages, there exist comprehensive and clear style guides.%
\item<2-> Such style guides tell us how to name things and how to structure code consistently.%
\item<3-> This is important, because we always work together with other people on projects and everybody should be able to read everyone's code.%
\item<4-> For SQL, there does not exist one generally accepted best practice standard style guide.%
\item<5-> We will still try to define some general rules.%
\end{itemize}%
\uncover<6->{%
\bestPractice{tableName}{%
Tabellennamen sollten (Englische) Nomen im Singular sein, mit Kleinbuchstaben geschrieben, und ohne Prefix~(z.B.~kein \inQuotes{tbl\_} am Anfange haben)~\cite{B2025DS:SBPASG}.%
}%
\uncover<7->{\begin{itemize}%
\item Nennen wir unsere Tabelle für Produkte also \textil{products}.%
\end{itemize}%
}}%
\end{frame}%
%
\begin{frame}[t]%
\frametitle{Grobstruktur}%
\begin{itemize}%
\only<-10>{%
\item Erstellen wir also die Tabelle~\sqlil{product}.%
\item<2-> Welche Spalten brauchen wir?%
\uncover<3->{%
\begin{itemize}%
\item Jedes Produkt hat einen \emph{Namen}.\uncover<4->{ %
Der Name ist im Grunde Text, also brauchen wir eine Spalte für Text.}%
\item<5-> Jedes Produkt hat einen Preis.\uncover<6->{ %
Der Preis ist eine Zahl die den Anforderungen an eine Geldmenge entsprechen muss.}%
\item<7-> Um die Produkte auszuliefern, kommt jedes Produkt in ein Schachtel.\uncover<8->{ %
Diese Schachteln haben eine Breite, Höhe, und Tiefe, sowie (wenn das Produkt eingepackt ist) ein bestimmtes Gewicht.\uncover<9->{ %
Das sind also noch vier Zahlen, die wir speichern müssen.}}%
\end{itemize}%
}}%
\item<10-> Das ist also die grundlegende Struktur unserer Tabelle.%
\end{itemize}%
\locateGraphic{11}{width=0.625\paperwidth}{graphics/create/dbStructureProduct}{0.1875}{0.2}%
\end{frame}%
%
\begin{frame}%
\frametitle{SQL}%
\begin{itemize}%
\item Diese Struktur ist im Grunde recht einfach.%
\item<2-> Nun müssen wir sie in SQL-Kommandos übersetzen.%
\item<3-> Wir müssen dazu Namen, Datentypen, und Einschränkungen für die Spalten festlegen.%
\item<4-> Dann können wir die SQL-Kommandos an das \postgresql\ DBMS schicken und die Tabelle erstellen.%
\item<5-> Dies kann dann innerhalb einer interaktiven \psql-Session oder mit Hilfe eines Skripts erfolgen.%
\item<6-> Schauen wir uns das mal an.%
\end{itemize}%
\end{frame}%
%
\begin{frame}[t]%
\frametitle{Interaktive PSQL Session -- Teil 1}%
\begin{itemize}%
\only<-2>{\item Wir erstellen das Kommando zuerst insgesamt und erklären dann Schritt-für-Schritt seine Komponenten.%
\item<2-> Wir loggen uns jetzt in unseren \postgresql-Server ein und erstellen die neue Tabelle.}%
\only<3>{\item<3-> Wir starten \psql\ mit der passenden Verbindungs-URI. Diesmal verwenden with den Benutzer \textil{boss} mit Passwort \textil{superboss123} und arbeiten auf der Datenbank \textil{factory}.}%
\only<4>{\item<4-> \psql\ läuft und ist mit dem \postgresql\ DBMS verbunden.}%
\only<5>{\item<5-> Mit diesem \sqlil{CREATE TABLE} Kommando -- das wir gleich im Detail anschauen -- wird die Tabelle mit allen notwendigen Spalten erstellt.}%
\only<6>{\item<6-> Das Kommando wird erfolgreich ausgeführt und \sqlil{CREATE TABLE} wird uns als Antwort ausgegeben.}%
\end{itemize}%
%
\locateGraphic{3}{width=0.85\paperwidth}{graphics/create/createTable1psql}{0.075}{0.25}%
\locateGraphic{4}{width=0.85\paperwidth}{graphics/create/createTable2started}{0.075}{0.25}%
\locateGraphic{5}{width=0.85\paperwidth}{graphics/create/createTable3command}{0.075}{0.25}%
\locateGraphic{6}{width=0.85\paperwidth}{graphics/create/createTable4created}{0.075}{0.25}%
\end{frame}%
%
\begin{frame}[t]%
\frametitle{Command Structure}%
\begin{itemize}%
\only<-1>{\item \sqlil{CREATE TABLE} ist das SQL-Kommando um eine Tabelle anzulegen.}%
\only<2>{\item<2> \sqlil{product} ist der Name der Tabelle, die wir anlegen wollen.}%
\only<3>{\item<3> Dann folgen die Namen, Datentypen, und Einschränkungen der Spalten unserer Tabelle. %
Diese sind von Klammern umrahmt~(\textil{(...)}) und durch Kommas~(\textil{,}) getrennt.}%
\end{itemize}%
\locateGraphic{1}{width=0.85\paperwidth}{graphics/create/createTable1psql}{0.075}{0.275}%
\locateGraphic{2}{width=0.85\paperwidth}{graphics/create/createTable3commandTableName}{0.075}{0.275}%
\locateGraphic{3}{width=0.85\paperwidth}{graphics/create/createTable3commandSpalten}{0.075}{0.275}%
\end{frame}%
%
\begin{frame}[t]%
\frametitle{Die Spalte \texttt{name}}%
\begin{itemize}%
%
\only<-4>{\item Wir legen eine Spalte \textil{name} für die Namen der Produkte an.%
\uncover<2->{ Die Namen sind Text, haben aber keine feste Lände.%
\uncover<3->{ Dafür gibt es den Datentyp \sqlil{VARCHAR}.%
\uncover<4->{ In Klammern geben wir eine Maximallänge an. Hier ist 100 Zeichen eine vernünftige Wahl.%
}}}}%
%
\only<5-8>{\item<5-> Jedes Produkt muss einen Namen haben.%
\uncover<6->{ Es kann kein Produkt ohne Name geben.%
\uncover<7->{ Wir schreiben daher \sqlil{NOT NULL} hinter den Datentypen.%
\uncover<8->{ Diese Einschränkung stellt sicher, dass jeder Datensatz immer einen Wert in der Spalte \sqlil{name} haben muss.%
}}}}%
%
\only<9-12>{\item<5-> Die Produktnamen müssen eindeutig sein.%
\uncover<10->{ Zwei verschiedene Produkte dürfen niemals den gleichen Namen haben.%
\uncover<11->{ Wir schreiben daher \sqlil{UNIQUE} hinter den Datentypen.%
\uncover<12->{ Diese Einschränkung stellt sicher, dass jeder Wert in der Spalte \sqlil{name} nur einmal vorkommen kann.%
}}}}%
%
\end{itemize}%
\locateGraphic{-4}{width=0.8\paperwidth}{graphics/create/createTable3commandNameVarchar}{0.1}{0.33}%
\locateGraphic{5-8}{width=0.8\paperwidth}{graphics/create/createTable3commandNameNotNull}{0.1}{0.33}%
\locateGraphic{9-12}{width=0.8\paperwidth}{graphics/create/createTable3commandNameUnique}{0.1}{0.33}%
\end{frame}%
%
\begin{frame}[t]%
\frametitle{Die Spalte \texttt{price}}%
\begin{itemize}%
%
\only<-17>{\item Wir legen eine Spalte \textil{price} für den Preis der Produkte an.%
\uncover<2->{ Die Auswahl des richtigen Datentypen ist diesmal schwer. %
\only<3-15>{\begin{itemize}%
\item Preise wie \inQuotes{1,34€} oder \inQuotes{12.99元} oder \inQuotes{\$3.99} sehen auf den ersten Blick wie Fließkommazahlen aus.\only<4-7>{%
\begin{itemize}%
\item Sie kennen Fließkommazahlen aus \python~(Datentype \pythonil{float}) oder \texttt{Java} und \textil{C}~(Datentypen \pythonil{float}, \pythonil{double}, \dots).%
\item<5-> Aber Fließkommazahlen sind ungenau!%
\item<6-> Beispiel:~\pythonil{0.1 + 0.1 + 0.1 - 0.3} ergiebt~\pythonil{5.551115123125783e-17} in \python\dots%
\end{itemize}%
}%
\only<8->{ Preise sind aber keine Fließkommazahlen!}%
%
\item<9-> Preise könnten Ganzzahlen sein, z.B. die Anzahl Cent oder 分.\only<10-11>{%
\begin{itemize}%
\item Dann muss man sie aber jedesmal in \inQuotes{umrechnen}, wenn man sie anzeigen oder mit ihnen rechnen will.%
\item<11-> Aber das ist fehleranfällig\cite{W2020HSISCVISS}.%
\end{itemize}%
}%
\only<12->{ Preise sind aber keine Ganzzahlen.}%
%
\only<13-15>{\item Der Datentype \sqlil{DECIMAL(x, y)} erlaubt es uns, Zahlen mit \textil{x}~Stellen (\textil{y} davon nach dem Komma) exakt darzustellen.%
\only<15>{%
\begin{itemize}%
\item Mit \sqlil{price DECIMAL(10, 2)} können wir Zahlen mit 10~Ziffern, 2~davon nach dem Komma, exakt speichern\cite{PGDG:PD:NT}.%
\end{itemize}%
}%
}%
%
\end{itemize}%
%
}%
\uncover<16->{ Wir benutzen den Datentyp~\sqlil{DECIMAL(10, 2)}.}%
}}%
%
\only<18-21>{\item Jedes Produkt muss einen Preis haben.%
\uncover<19->{ Es kann kein Produkt ohne Preis geben.%
\uncover<20->{ Wir schreiben daher \sqlil{NOT NULL} hinter den Datentypen.%
\uncover<21->{ Diese Einschränkung stellt sicher, dass jeder Datensatz immer einen Wert in der Spalte \sqlil{preis} haben muss.%
}}}}%
%
\end{itemize}%
%
\only<5-6>{%
\bestPractice{floatImprecise}{%
Nehmen Sie immer an, dass \pythonilIdx{float}-Werte ungenau sind. %
Nehmen Sie NIEMALS an, dass diese exact sind\cite{PTVF2007NRTAOSC:EAAS,BHK2006CNFCMEARFCS:NS,programmingWithPython}.%
}}%
\only<7-8>{%
\bestPractice{moneyNotFloat}{%
Repräsentieren Sie niemals Geldwerte mit Fließkommazahlen\cite{SE:SO:WNUDOFTRC,W2020HSISCVISS}.%
}}%
\only<14-16>{%
\bestPractice{moneyDecimal}{%
Geldwerte müssen mit dem Datentype \sqlilIdx{DECIMAL} gespeicher werden\cite{C2024SCVDTCBP,W2020HSISCVISS}.%
}}%
%
\locateGraphic{-3,17}{width=0.8\paperwidth}{graphics/create/createTable3commandPriceDecimal}{0.1}{0.343}%
\locateGraphic{18-21}{width=0.8\paperwidth}{graphics/create/createTable3commandPriceNotNull}{0.1}{0.343}%
\end{frame}%
%
%
\begin{frame}[t]%
\frametitle{Die Spalten \texttt{weight}, \texttt{width}, \texttt{height}, und \texttt{depth}}%
\begin{itemize}%
%
\only<-8>{\item Wir legen die Spalten für das Gewicht und die Dimensionen der Boxen für unsere Produkte an.%
\only<2-5>{%
\begin{itemize}%
\item Das Gewicht soll in Gram in der Spalte \textil{weight} gespeichert werden.%
\item<3-> Die Breite soll in Millimeter in der Spalte \textil{width} gespeichert werden.%
\item<4-> Die Höhe soll in Millimeter in der Spalte \textil{height} gespeichert werden.%
\item<5-> Die Tiefe soll in Millimeter in der Spalte \textil{depth} gespeichert werden.%
\end{itemize}%
\uncover<6->{ Alle Dimensionen sind Ganzzahlen.
\uncover<7->{ Der Datentype \sqlil{INT} stellt Werte von -2\decSep147\decSep483\decSep648 bis +2\decSep147\decSep483\decSep647 dar.
\uncover<8->{ Das reicht.}}}%
}%
}%
%
\only<9-12>{\item Jedes Produkt muss ein Gewicht und alle Box-Dimensionen haben.%
\uncover<10->{ Es kann kein Produkt ohne Gewicht oder mit unbekannter Größe geben.%
\uncover<11->{ Wir schreiben daher \sqlil{NOT NULL} hinter die Datentypen.%
\uncover<12->{ Diese Einschränkung stellt sicher, dass jeder Datensatz immer Werte für alle vier Attribute haben muss.%
}}}}%
%
\end{itemize}%
\locateGraphic{1,6-8}{width=0.8\paperwidth}{graphics/create/createTable3commandWWHDint}{0.1}{0.33}%
\locateGraphic{9-12}{width=0.8\paperwidth}{graphics/create/createTable3commandWWHDnotNull}{0.1}{0.355}%
\end{frame}%
%
%
\begin{frame}[t]%
\frametitle{Primärschlüssel}%
\begin{itemize}%
\item Die Tabellen in einer Datenbank existieren nicht alleine oder losgelöst.%
\item<2-> Sie referenzieren einander.%
\item<3-> Später wollen wir speichern, welcher Kunde welches Produkt bestellt hat.%
\item<4-> Dazu brauchen wir mindestens drei Tabellen\uncover<5->{:%
\begin{itemize}%
\item eine Tabelle für Kunden (die wir noch nicht haben),%
\item<6-> die Tabelle für Produkte (die wir gerade erstellen),%
\item<7-> und eine Tabelle, die eine Zeile aus der Kundentabelle mit einer Zeile aus der Produkttabelle zu einer Bestellung \inQuotes{verbindet}.%
\end{itemize}%
}%
\item<8-> Damit das geht, müssen wir die Zeilen in der Kunden- und der Produkttabelle finden können.%
\item<9-> Jede Zeile muss eindeutig identifiziert werden können.%
\item<10-> Dafür gibt es \emph{Primärschlüssel}.%
\end{itemize}%%
\end{frame}%
%
%
\begin{frame}[t]%
\frametitle{Primärschlüssel}%
\begin{itemize}%
\item Primärschlüssel sind (eine oder mehrere) Spalten einer Tabelle deren Werte die Zeilen eindeutig identifizieren.%
\item<2-> Das heißt, dass es keine zwei Zeilen einer Tabelle mit den selben Primärschlüsselwerten geben kann.%
\item<3-> Wir haben bereits so eine Spalte erstellt\uncover<4->{: \sqlil{name}}!%
\item<4-> Die Produktnamen, gespeichert in \sqlil{name}, sind einmalig~(\sqlil{UNIQUE}).%
\item<5-> Wir könnten diese Spalte als Primärschlüssel nehmen.%
\item<6-> Dann würde die Tabelle für Bestellungen für jede Bestellung den Namen des Produktes mitspeichern.%
\item<7-> Über den Namen können wir dann die Eigenschaften, z.B.\ den Preis, auslesen.%
\item<8-> Aber was, wenn wir irgendwann den Namen eines Produktes ändern?%
\item<9-> Dann geht die Verbindung verloren\dots%
\end{itemize}%
\end{frame}%
%
%
\begin{frame}[t]%
\frametitle{Die Spalte \texttt{id}}%
\begin{itemize}%
%
\only<-4>{\item Wir legen eine zusätzliche Spalte \textil{id} als Primärschlüssel an.%
\uncover<2->{ Wir markieren die Spalte as \sqlil{PRIMARY KEY}.%
\uncover<3->{ Dies impliziert \sqlil{NOT NULL} und \sqlil{UNIQUE}.%
\uncover<4->{ Aber was kommt in diese Spalte?%
}}}}%
%
\only<5-8>{\item<5-> Geben wir ihr ebenfalls den Datentyp \sqlil{INT}.%
\uncover<6->{ Die Primärschlüssel sollen also ganzzahlige Werte sein.%
\uncover<7->{ Anders als die Produktnamen, werden wir diese niemals ändern.%
\uncover<8->{ Aber was kommt in diese Spalte?%
}}}}%
%
\only<9-12>{\item<5-> Wir lassen sie automatisch von der Datenbank generieren!%
\uncover<10->{ \sqlil{GENERATED BY DEFAULT AS IDENTITY} sagt dem DBMS, dass wir diese Werte nicht selber angeben wollen\dots%
\uncover<11->{ {\dots}stattdessen soll es selbstständig einzigartige Werte eintragen.%
\uncover<12->{ Wir brauchen uns also nicht weiter darum zu kümmern.%
}}}}%
%
\end{itemize}%
\locateGraphic{-4}{width=0.8\paperwidth}{graphics/create/createTable3commandIdPK}{0.1}{0.35}%
\locateGraphic{5-8}{width=0.8\paperwidth}{graphics/create/createTable3commandIdInt}{0.1}{0.35}%
\locateGraphic{9-12}{width=0.8\paperwidth}{graphics/create/createTable3commandIdPKgenerated}{0.1}{0.35}%
\end{frame}%
%
%
\begin{frame}[t]%
\frametitle{Interaktive PSQL Session -- Teil 2}%
\begin{itemize}%
\only<-1>{\item Damit ist das SQL-Kommando erklärt.}%
\only<2>{\item<2-> Wie bereits gesagt können wir es erfolgreich auf unserem \postgresql-Server ausführen.}%
\only<3-4>{\item<3-> Mit einer \sqlil{SELECT ... FROM} Anfrage auf die passende Systemtabelle können wir prüfen, ob die Tabelle \sqlil{product} nun wirklich existiert.%
\uncover<4->{ Sie tut es.}}%
\only<5>{\item<5-> Wir beenden die \psql-Session durch \textil{\\q} und \keys{\enter}.}%
\only<6>{\item<6-> Wir sind zurück im normalen Terminal.}%
\end{itemize}%
%
\locateGraphic{-2}{width=0.85\paperwidth}{graphics/create/createTable4created}{0.075}{0.25}%
\locateGraphic{3}{width=0.85\paperwidth}{graphics/create/createTable5select}{0.075}{0.25}%
\locateGraphic{4}{width=0.85\paperwidth}{graphics/create/createTable6selected}{0.075}{0.25}%
\locateGraphic{5}{width=0.85\paperwidth}{graphics/create/createTable7quit}{0.075}{0.25}%
\locateGraphic{6}{width=0.85\paperwidth}{graphics/create/createTable8done}{0.075}{0.25}%
\end{frame}%
%
%
\begin{frame}[fragile]
\frametitle{Erstellen der Tabelle via Script}%
\gitLoadAndExecSQL{factory:create_table_product}{}{factory}{create_table_product.sql}{factory}{boss}{superboss123}%
\listingSQLandOutput{1}{factory:create_table_product}{}{0.225}{0.07}{0.9}{0.91}
\listingSQL{2}{factory:create_table_product}{0.05}{0.07}{0.9}{0.91}
\listingOutput{3}{factory:create_table_product}{,style=tool_style_sql}{0.05}{0.07}{0.9}{0.91}
\end{frame}
%
\section{Einfügen von Daten}%
%
\begin{frame}%
\frametitle{Daten Einfügen}%
\begin{itemize}%
\item Nun haben wir die Tabelle \sqlil{product} erstellt, aber sie ist leer.%
\item<2-> Füllen wir sie mit Daten!%
\item<3-> Unsere Fabrik hat zwei Produkte:~Schuhe~(Shoe) und Handtaschen~(Purse).%
\item<4-> Schuhe\uncover<4->{:%
\begin{itemize}%
\item Schuhe bieten wir in den Schuhgrößen~36 bis~43 an.%
\item<5-> Ihr Preis beginnt bei 150.99元 für Schuhgröße~36 und steigt um 2元 pro Größe.%
\item<6-> Die Schuhe passen alle in diese selbe Box.%
\item<7-> Die kleinsten Schuhe wiegen 1300g und das Gewicht steigt um 25g pro Größe.%
\end{itemize}}%
%
\item<8-> Handtaschen\uncover<8->{:%
\begin{itemize}%
\item Handtaschen kommen in den Größen klein (\emph{small}), mittel (\emph{medium}), und groß (\emph{large}).%
\item<9-> Die entsprechenden Preise sind 100元, 120元, und 150元.%
\item<10-> Die kleinste Handtasche passt in einen Schuhkarton, die größeren erfordern größere Boxen.%
\item<11-> Die entsprechenden Gewichte sind dann 500g, 750g, und 1500g.%
\end{itemize}%
}%
\end{itemize}%
\end{frame}%
%
%
\section{Zusammenfassung}%
%
\begin{frame}%
\frametitle{Zusammenfassung}%
\begin{itemize}%
\item Fertig
\end{itemize}%
\end{frame}%
%
\endPresentation%
\end{document}%%
\endinput%
%
